\section{Posiciones de Bloques}

\lead{La protección de la infancia en el entorno digital trasciende fronteras nacionales y requiere la participación coordinada de todos los involucrados.}

Sin embargo, los Estados mantienen posiciones divergentes respecto al alcance de la regulación, el equilibrio entre seguridad y derechos fundamentales y el papel que deben desempeñar las plataformas digitales. Estas diferencias corresponden a factores políticos, jurídicos, económicos y tecnológicos.

\subsection{Estados con enfoques regulatorios estrictos: Unión Europea}

La Unión Europea (UE) ha adoptado uno de los enfoques más intervencionistas en materia de regulación digital, priorizando la protección de los menores y la responsabilidad de las plataformas tecnológicas. Su enfoque se basa en la responsabilidad proactiva: las empresas deben demostrar que sus espacios son seguros antes de lanzarlos al mercado.

La entrada en vigor del DSA y el GDPR refleja la filosofía sobre la que la UE se basa.\footnote{European Commission. (2024). Reglamento de Servicios Digitales (Digital Services Act). \url{https://digital-strategy.ec.europa.eu/en/policies/digital-services-act}} Estas normas obligan a las grandes plataformas a evaluar y mitigar riesgos sistémicos, fortalecer los mecanismos de denuncia de contenido ilegal, aumentar la transparencia de sus algoritmos y proteger específicamente a los menores frente a contenidos perjudiciales.

Cabe resaltar que, aunque se está hablando específicamente de miembros de la Unión Europea, estas no son las únicas naciones que tienen este enfoque.

Otros países como el Reino Unido o Australia también han demostrado participación activa en la protección de los menores antes que explotar los intereses comerciales de las plataformas.

\subsection{Estados que priorizan la privacidad y la libertad de expresión: Estados Unidos}

Este bloque comparte la urgencia de proteger a la infancia, pero difiere radicalmente en los métodos. Estados Unidos es el mayor ejemplo de esto. Su posición defiende la Primera Enmienda y los derechos de privacidad globales, oponiéndose a legislaciones que obliguen a las plataformas a escanear de forma masiva los mensajes privados de los ciudadanos o que restrinjan el cifrado de extremo a extremo, argumentando que esto debilitaría la seguridad cibernética global y abriría la puerta a espionaje estatal.

También, al ser Estados Unidos albergar la sede de muchas de las principales empresas tecnológicas del mundo, el gobierno estadounidense suele favorecer la innovación tecnológica y la protección de la libertad de expresión, procurando evitar regulaciones que puedan interferir con el desarrollo de estas compañías. No obstante, el país cuenta con una amplia experiencia en la persecución de delitos de explotación infantil mediante organismos como el National Center for Missing and Exploited Children (NCMEC) y el Federal Bureau of Investigation (FBI), que colaboran estrechamente con empresas tecnológicas y cuerpos policiales internacionales para identificar víctimas y desmantelar redes criminales.\footnote{Missing \& Exploited Children. (s. f.). Datos de la CyberTipline. Centro Nacional para Niños Desaparecidos y Explotados. \url{https://www.missingkids.org/gethelpnow/cybertipline/cybertiplinedata}}

\subsection{Estados con un enfoque de fuerte control estatal: China y otros modelos similares}

Algunos países, como China, adoptan una visión donde la seguridad nacional y la estabilidad social justifican una intervención mucho más amplia del Estado en el entorno digital. Estos gobiernos consideran que la regulación estricta de internet constituye una herramienta legítima para prevenir tanto la radicalización como la explotación infantil.

China mantiene uno de los sistemas de supervisión digital más extensos del mundo, con controles sobre plataformas nacionales, restricciones al acceso de determinadas aplicaciones extranjeras y amplias facultades para monitorear contenidos considerados ilegales o perjudiciales.

En relación a la protección infantil, el gobierno ha implementado medidas como restricciones al tiempo de uso de videojuegos para menores, sistemas de verificación de identidad y obligaciones estrictas para empresas tecnológicas respecto a la moderación de contenidos.

Aunque estas políticas son eficaces para proteger a los menores, diversos organismos internacionales y organizaciones defensoras de los derechos humanos han expresado preocupación por el posible impacto de estas medidas sobre la privacidad, la libertad de expresión y el acceso a la información.

\subsection{Países en desarrollo}

Los países en desarrollo y economías emergentes se enfrentan a una dura realidad; la falta de recursos institucionales, económicos y tecnológicos para examinar e inspeccionar en esta área. Muchos países de África, América Latina, el Caribe y partes de Asia reconocen la gravedad del problema y su prioridad es cerrar la brecha digital y fomentar el crecimiento económico a través de la conectividad, por lo que las regulaciones extremadamente complejas o costosas pueden aislar tecnológicamente a sus ciudadanos. Es por este motivo que estos estados suelen enfatizar la necesidad de cooperación internacional, transferencia de tecnología, asistencia financiera y fortalecimiento institucional.\footnote{World Economic Forum. (2025, enero). Avances tecnológicos y desarrollo humano: Una historia de dos mundos. \url{https://es.weforum.org/stories/2025/01/avances-tecnologicos-y-desarrollo-humano-una-historia-de-dos-m} undos/}

Del mismo modo, muchos de estos países consideran prioritario que las empresas tecnológicas asuman una mayor responsabilidad, dado que gran parte de las plataformas utilizadas por sus ciudadanos tienen sede en países desarrollados.
